\chapter[Mission, Project Management and Personal Contribution]
{Mission, Project Management and\\Personal Contribution}
\label{chap:mission}

\section{Mission and work packages}

The internship title used in this report is
\emph{\ReportMissionTitle}. The signed agreement identifies the placement,
host, dates and entrusted activities \citep{internshipagreement}. The work
reported here concerns simulation; no deployment on physical platforms is
claimed. The first practical assignment was to model Room~315 in Gazebo and
establish a usable simulation environment for its robotic and transport
equipment. Completing this objective ahead of schedule left time for a research
work package combining overhead-camera perception with natural-language task
interaction. The internship dates were unchanged.

The implemented work is organised into two phases:

\begin{description}
  \item[Phase 1---simulation foundation] Room~315 Gazebo modelling and a
  reusable, selectively launched multi-robot and rail simulation foundation.
  \item[Phase 2---neuro-symbolic research] Paired cameras, visual-state
  estimation, validated language goals, PDDL/PlanSys2 and safety-gated
  neuro-symbolic control, in which learned proposals are checked before
  symbolic planning and software-supervised execution.
\end{description}

The phases describe the project chronology, whereas objectives O1--O4 in
\cref{chap:introduction} define the units used to assess the completed work.
The simulation phase established the multi-robot and transport foundations.
The neuro-symbolic phase extended the typed interfaces and
connected visual and language inputs to planning, software-supervised execution and
effect verification.

\section{Chronology and decision milestones}

The internship schedule was fixed from the outset. As shown in
\cref{fig:timeline}, the technical phases overlap because integration, testing
and correction continued after each main decision milestone.

\begin{figure}[htbp]
  \centering
  \resizebox{\textwidth}{!}{\begin{tikzpicture}[
  font=\sffamily\footnotesize,
  panel/.style={draw=mfjablue!35,fill=white,rounded corners=3pt,
    line width=0.6pt,inner sep=0pt},
  recorded/.style={draw=mfjablue,fill=mfjablue!17,rounded corners=1.2pt,
    line width=0.7pt},
  planned/.style={draw=mfjaorange!90!black,fill=mfjaorange!6,
    rounded corners=1.2pt,line width=0.8pt,densely dashed},
  gridline/.style={draw=mfjablue!15,line width=0.4pt},
  milestone/.style={draw=mfjagreen!75!black,line width=0.65pt,
    densely dashed},
  statusline/.style={draw=black!55,line width=0.75pt}
]

% A single calendar scale is used throughout.  Horizontal positions are
% proportional to the 179 days from 2 March (x=-20 mm) to 28 August (x=74 mm).
\node[panel,minimum width=160mm,minimum height=130mm] (timeline) at (0,-5mm) {};

\node[anchor=west,font=\sffamily\bfseries\large,text=mfjablue]
  at (-75mm,54mm) {Internship chronology};
\node[anchor=east,font=\sffamily\small,text=black!62]
  at (75mm,54mm) {Status at 18 August 2026};

% Shade the interval after the report-review date before drawing the grid.
\path[fill=mfjaorange!3] (68.75mm,30mm) rectangle (74mm,-58mm);

% Common calendar axis and month grid.
\draw[mfjablue,line width=0.8pt] (-20mm,35mm) -- (74mm,35mm);
\foreach \x/\month in {
  -20/{2 Mar},
  -4.25/{Apr},
  11.51/{May},
  27.79/{Jun},
  43.55/{Jul},
  59.82/{Aug},
  74/{28 Aug}}
{
  \draw[mfjablue,line width=0.65pt] (\x mm,33.8mm) -- (\x mm,36.2mm);
  \node[anchor=north,font=\sffamily\scriptsize,text=black!68]
    at (\x mm,32.8mm) {\month};
  \draw[gridline] (\x mm,30mm) -- (\x mm,-58mm);
}

% Evidence-backed decision milestones and the completed Final Test.
\draw[milestone] (-7.92mm,35mm) -- (-7.92mm,43.5mm);
\fill[mfjagreen!75!black] (-7.92mm,35mm) circle (1.05pt);
\node[anchor=south,align=center,font=\sffamily\scriptsize,
  text=mfjagreen!65!black] at (-7.92mm,44mm)
  {25 Mar\\package architecture};

\draw[milestone] (14.66mm,35mm) -- (14.66mm,39.5mm);
\fill[mfjagreen!75!black] (14.66mm,35mm) circle (1.05pt);
\node[anchor=south,align=center,font=\sffamily\scriptsize,
  text=mfjagreen!65!black] at (14.66mm,40mm)
  {7 May\\typed rail interfaces};

\draw[milestone] (48.79mm,35mm) -- (48.79mm,43.5mm);
\fill[mfjagreen!75!black] (48.79mm,35mm) circle (1.05pt);
\node[anchor=south east,align=right,font=\sffamily\scriptsize,
  text=mfjagreen!65!black] at (54.5mm,44mm)
  {11 Jul\\closed-loop contracts};

\draw[milestone] (65.07mm,35mm) -- (65.07mm,39.5mm);
\fill[mfjagreen!75!black] (65.07mm,35mm) circle (1.05pt);
\node[anchor=south west,align=left,font=\sffamily\scriptsize,
  text=mfjagreen!65!black] at (59.5mm,40mm)
  {11 Aug\\Final Test};

% Phase 1: platform construction and integration.
\node[anchor=east,align=right,font=\sffamily\scriptsize\bfseries,text=mfjablue]
  at (-23mm,27mm) {PHASE 1 --- PLATFORM \& INTEGRATION};
\foreach \y/\label in {
  21/{CAD/Gazebo scene and robot setup},
  15/{Multi-robot launch and isolation},
  9/{Rail topology, motion and devices},
  3/{Typed interfaces and safety checks},
  -3/{Camera data and scene estimation}}
{
  \node[anchor=east,align=right] at (-24mm,\y mm) {\label};
}
\path[recorded] (-20mm,19.1mm) rectangle (-7.92mm,22.9mm);
\path[recorded] (-18.42mm,13.1mm) rectangle (29.36mm,16.9mm);
\path[recorded] (-7.92mm,7.1mm) rectangle (29.36mm,10.9mm);
\path[recorded] (14.66mm,1.1mm) rectangle (51.95mm,4.9mm);
\path[recorded] (29.36mm,-4.9mm) rectangle (51.95mm,-1.1mm);

\draw[mfjablue!18,line width=0.45pt] (-75mm,-7.5mm) -- (74mm,-7.5mm);

% Phase 2: planning, validation and the fixed evidence gate.
\node[anchor=east,align=right,font=\sffamily\scriptsize\bfseries,text=mfjablue]
  at (-23mm,-11mm) {PHASE 2 --- PLANNING \& VALIDATION};
\node[anchor=east] at (-24mm,-17mm) {PDDL/PlanSys2 planning};
\node[anchor=east] at (-24mm,-23mm) {Validation and Final Test};
\path[recorded] (34.61mm,-18.9mm) rectangle (62.97mm,-15.1mm);
\path[recorded] (48.79mm,-24.9mm) rectangle (65.07mm,-21.1mm);

\draw[mfjablue!18,line width=0.45pt] (-75mm,-27.5mm) -- (74mm,-27.5mm);

% Reporting shares the same scale; the short final interval is therefore not
% exaggerated.  Only the period after 18 August is styled as planned work.
\node[anchor=east,align=right,font=\sffamily\scriptsize\bfseries,text=mfjablue]
  at (-23mm,-31mm) {DOCUMENTATION AND REPORTING};
\node[anchor=east] at (-24mm,-37mm)
  {Project documentation};
\node[anchor=east] at (-24mm,-43mm)
  {Report section drafting};
\node[anchor=east] at (-24mm,-49mm)
  {Final report assembly and revision};
\node[anchor=east] at (-24mm,-55mm)
  {Defence preparation};
\path[recorded] (-20mm,-38.9mm) rectangle (68.75mm,-35.1mm);
\path[recorded] (-20mm,-44.9mm) rectangle (61.40mm,-41.1mm);
\path[recorded] (61.40mm,-50.9mm) rectangle (68.75mm,-47.1mm);
\path[planned] (68.75mm,-56.9mm) rectangle (74mm,-53.1mm);

% The review boundary separates recorded work from the remaining plan.
\draw[statusline] (68.75mm,30mm) -- (68.75mm,-58mm);
\node[rotate=90,anchor=south,font=\sffamily\scriptsize\bfseries,
  text=black!62,fill=white,inner xsep=1.2pt,inner ysep=0.5pt]
  at (68.75mm,-14mm) {18 Aug review};

% Compact legend.
\path[recorded] (-73mm,-65mm) rectangle (-68.5mm,-61.6mm);
\node[anchor=west,font=\sffamily\scriptsize,text=black!68]
  at (-66.5mm,-63.3mm) {Recorded activity period};
\path[planned] (-13mm,-65mm) rectangle (-8.5mm,-61.6mm);
\node[anchor=west,font=\sffamily\scriptsize,text=black!68]
  at (-6.5mm,-63.3mm) {Planned after 18 Aug};
\fill[mfjagreen!75!black] (40mm,-63.3mm) circle (1.05pt);
\node[anchor=west,font=\sffamily\scriptsize,text=black!68]
  at (43mm,-63.3mm) {Milestone / evidence gate};

\end{tikzpicture}
}
  \caption[Engineering, reporting and defence-preparation chronology.]
  {Internship chronology as of 11 August 2026. Engineering, integration,
  visual-model validation and the preregistered Final Test extend through 11 August in
  panel~(a), overlapping the report-focused period that began on 4 August.
  Documentation continues to 18 August; the final period, 18--28 August, is
  reserved for defence-presentation preparation.}
  \label{fig:timeline}
\end{figure}

\section{Decision progression}

\Cref{fig:stage-decision-chain} summarises how the project progressed from the
initial simulation to the integrated neuro-symbolic system. Each stage depended
on working results from the preceding stage.

Major technical decisions were taken after weekly discussions with the
supervisors. The agreed solutions were then implemented and checked through
configuration, tests and integration runs.

\begin{figure}[htbp]
  \centering
  \resizebox{0.98\textwidth}{!}{\begin{tikzpicture}[
  font=\sffamily\small,
  stage/.style={
    draw=black!22,
    rounded corners=2mm,
    fill=white,
    align=left,
    text width=151mm,
    minimum height=18mm,
    inner sep=2.4mm,
    execute at begin node={\raggedright
      \hyphenpenalty=10000\exhyphenpenalty=10000
      \emergencystretch=3em\relax}
  },
  phase/.style={
    rounded corners=1.2mm,
    fill=mfjablue,
    text=white,
    font=\sffamily\bfseries,
    text width=151mm,
    align=left,
    inner xsep=2.4mm,
    inner ysep=1.2mm
  },
  flow/.style={-{Latex[length=2.2mm]},line width=1.05pt,mfjablue}
]
  \node[phase] (p1) {Phase 1 --- initial assigned task};

  \node[stage,below=3mm of p1] (s1) {%
    {\color{mfjablue}\bfseries 1. Third-floor CAD and robots}\par\smallskip
    {\color{mfjaorange}\bfseries Engineering need:} create the floor environment and give different robot models consistent names and communication channels.\quad
    {\color{mfjablue}\bfseries Implementation:} design the floor geometry and room furnishings in Autodesk Inventor, export them to Gazebo, separate software responsibilities
    and give each robot a separate set of names.\quad
    {\color{mfjagreen}\bfseries Result:} a furnished full-floor scene that can start any supported combination of simulated robots.};

  \node[stage,below=4mm of s1] (s2) {%
    {\color{mfjablue}\bfseries 2. Shuttle motion}\par\smallskip
    {\color{mfjaorange}\bfseries Engineering need:} repeatable motion through rail switches.\quad
    {\color{mfjablue}\bfseries Implementation:} use smooth rail paths parameterised by travelled distance and Gazebo commands for position and orientation.\quad
    {\color{mfjagreen}\bfseries Result:} repeatable continuous motion with an explicitly
    defined path-following model.};

  \node[stage,below=4mm of s2] (s3) {%
    {\color{mfjablue}\bfseries 3. Rail layout and devices}\par\smallskip
    {\color{mfjaorange}\bfseries Engineering need:} define allowed routes, stops and device
    behaviour.\quad
    {\color{mfjablue}\bfseries Implementation:} declare one-way rail segments, switches, slots, stoppers
    and sensors in configuration files.\quad
    {\color{mfjagreen}\bfseries Result:} state of both rails that the planner can read to coordinate shuttles.};

  \draw[flow] (s1.south) -- (s2.north);
  \draw[flow] (s2.south) -- (s3.north);

  \node[phase,below=5mm of s3] (p2)
    {Phase 2 --- additional research direction};

  \node[stage,below=3mm of p2] (s4) {%
    {\color{mfjablue}\bfseries 4. Camera and task interfaces}\par\smallskip
    {\color{mfjaorange}\bfseries Engineering need:} interpret English tasks and the current camera scene.\quad
    {\color{mfjablue}\bfseries Implementation:} pair overhead colour-camera views, keep reference labels separate from
    model input, and convert task descriptions with predefined fields into confirmed goals.\quad
    {\color{mfjagreen}\bfseries Result:} structured language and camera inputs for planning.};

  \node[stage,below=4mm of s4] (s5) {%
    {\color{mfjablue}\bfseries 5. Planning and software-supervised execution}\par\smallskip
    {\color{mfjaorange}\bfseries Engineering need:} convert the current task and equipment state into checked motion.\quad
    {\color{mfjablue}\bfseries Implementation:} rebuild the PDDL task description from the latest accepted state, select one named action,
    translate it into a command, pass it through the software supervisor and check the resulting state.\quad
    {\color{mfjagreen}\bfseries Result:} planning and executing one action at a time until completion.};

  \node[stage,below=4mm of s5] (s6) {%
    {\color{mfjablue}\bfseries 6. Experiments and results}\par\smallskip
    {\color{mfjaorange}\bfseries Engineering need:} measure scene estimation and confirm complete task behaviour.\quad
    {\color{mfjablue}\bfseries Implementation:} use predefined image-data groups, evaluation measures, software
    tests and 12 independent complete-task runs from a fresh start.\quad
    {\color{mfjagreen}\bfseries Result:} recorded model performance and completed runs for valid tasks across both
    rails, all eight shuttles and five task categories.};

  \draw[flow] (s4.south) -- (s5.north);
  \draw[flow] (s5.south) -- (s6.north);
\end{tikzpicture}
}
  \caption[Dependency-driven progression of the two project phases.]
  {Project progression as a dependency chain. Each stage converts a
  technical objective into an engineering result required by
  the next stage.}
  \label{fig:stage-decision-chain}
\end{figure}
\FloatBarrier

\section{Requirements and review process}

Requirements were expressed at three levels. Launch arguments, topics and
runbooks define how operators use the system. Executable contracts and tests
specify what each component must accept, reject or preserve. Artefact manifests
and hashes give datasets and models reproducible identities. The main
requirements across these levels were robot isolation without unintended
cross-namespace communication, stable shuttle identity across interfaces,
partial device updates without unintended changes, current state before planning,
deterministic translation from each symbolic plan step to a typed command and a
post-command observation after each executed plan step.

Each property was checked where it was implemented: geometry and coordinate
frames, interface semantics, shared-rail coordination and learned-model output.
The safeguards and data-leakage controls are detailed in
\cref{chap:transport,chap:interfaces,chap:ai-development}.

The agreed solutions were reflected in configuration and version-controlled
code, then exercised through the corresponding tests and experiments. The
identifiers for the relevant sources, datasets and models are listed in
\cref{app:reproducibility}; the corresponding objective-to-result assessment
is reported in \cref{tab:objective-results}.

\section{Personal responsibility, collaboration and resources}

Within the shared repository, I was responsible for the third-floor and
Room~315 models; the selective multi-robot launch architecture; calibrated rail
transport, device logic and typed interfaces; task contracts, PDDL/PlanSys2
integration and software-supervised one-step execution; the paired-camera dataset and
visual training pipeline; and the tests, manifests, runbooks and handover
documentation.

The supervisors provided scientific direction, technical review and knowledge
of the platform. During the weekly discussions, I presented the alternatives,
the available data and the consequences for shared interfaces, software
safeguards and experimental claims. After we agreed on a decision, I implemented it in
the relevant configuration, code, tests and documentation.

In practice, resource management concerned the fixed schedule, shared platform
access, simulation and training time, dataset and checkpoint storage, and the
testing workload. I kept compute-intensive perception and language nodes
optional, prioritised reusable configuration and regression tests, and prepared
build and run instructions so that the host could continue using the work.
No purchasing budget was assigned to me. The resources I had to manage were GPU
time, storage capacity, simulation time and access to the shared platform.

\section{Autonomy and reflective engineering examples}

I divided the broad integration objective into testable increments and kept
early implementation choices reversible. Before keeping a design, I checked
that its interfaces and implementation were consistent and that the corresponding
regression tests passed. I discussed decisions affecting shared interfaces,
experimental claims or software safeguards with the supervisors before
finalising them. The following two examples show how this worked in practice.

\subsection{Selecting a defensible level of rail fidelity}

At the beginning of the rail work, I compared a dynamic wheel--rail simulation
with a kinematic rail-following model. The measured mechanical parameters
required for a credible dynamic model were not available, and that option was
expected to require more integration time and computation. The project
questions instead concerned routing, sensor crossings, slot arrival, perception
and supervised execution. I therefore proposed the calibrated graph and
arc-length solution detailed in \cref{chap:transport}. I discussed the
comparison with the supervisors during a weekly review. The agreed priority was
repeatable behaviour at the level required by those questions. Mechanical
effects that could not be supported by the available data were left outside
the model. I then implemented
the model, checked the configured geometry and tested routing, sensor feedback
and slot-arrival behaviour. This experience showed me that model fidelity had
to follow the question being studied. In this case, the missing mechanical data,
the schedule and the available computing capacity were practical design
constraints.

\subsection{Choosing PDDL/PlanSys2 over end-to-end VLA control}

For the second phase, I compared an end-to-end vision--language--action model
that would infer actions directly with VLA and hybrid planning architectures
discussed with my supervisors. This comparison supported a clear division of
responsibilities: the visual model estimates the scene
state from the paired camera images, while PDDL/PlanSys2 determines the action
sequence, whose actions are then authorised by the software execution supervisor
and applied by the controllers. This
approach was better suited to Room~315 because its rail topology and device
rules are explicit, whereas the available dataset contains scene-state labels
rather than action demonstrations. It also made the plans easier to inspect and
execution errors easier to locate. This became the architecture implemented in
\cref{chap:ai-development,chap:neurosymbolic}.

\section{Competencies and professional learning}

\Cref{tab:competencies} relates the principal competencies developed during the
internship to concrete examples from the project.

\begin{table}[H]
\centering
\small
\caption{Engineering competencies and concrete project examples.}
\label{tab:competencies}
\begin{tabularx}{\textwidth}{@{}p{0.24\textwidth}X@{}}
\toprule
\textbf{Competency} & \textbf{Project example}\\
\midrule
Systems analysis & Separation of geometry, perception, planning, supervision
and exclusive command-publication authority\\
\tablerowrule
Robotics integration & ROS~2/Gazebo assets, controllers, launch composition,
namespaces and device feedback\\
\tablerowrule
Algorithm design & Arc-length graph, topology-aware routing, reservations,
clearance and effect checks\\
\tablerowrule
AI/ML method & Local language-model integration, leakage-controlled
paired-image data, declared partitions and separate visual-output heads\\
\tablerowrule
Software quality & Typed contracts, focused regression tests, configuration
manifests and reproducible result records\\
\tablerowrule
Project management and communication & Incremental milestones, supervisory
reviews, diagrams, runbooks, result tables and explicit scope\\
\bottomrule
\end{tabularx}
\end{table}

These two decisions changed how I compared technical alternatives. I used three
questions throughout the later integration work: what data were available, what
behaviour Room~315 required, and how the result could be tested.
